\section{Expensive layers, rank growth, and the contradiction}
\label{c:rank-section}

The bounded-width coordinates now have value zero. We extend their
independent initials through the remaining width layers, while keeping
discrepancy small and making the new coordinate values integral. Once
there are $d$ independent initials, a section count bounds the denominator
of the entire value group. This produces the prime divisor that contradicts
$\epsilon$-log canonicity. The
\hyperref[c:choice-order]{completion proof at the end of this section}
chooses the tolerances and the final small scaling after all uniform size
bounds have been established. Proposition~\ref{c:divisorial-index} uses a
separate large-degree limit on one fixed member of the sequence.

After reordering and passing to a subsequence, partition the
unbounded widths \(W_j\) into successive layers.  The widths in
one layer are comparable to a scale \(S\to\infty\), and each
scale divided by the next tends to zero.  All cheap widths are
comparable to \(1\).  If every width is cheap, the valuation
\(w_c\) constructed in the preceding section is Gaussian of
value zero on a full transcendence basis.  It is therefore
trivial on its rational field and on the algebraic extension
\(K\), contradicting its positive barrier on \(\m_z\).
We may assume that an expensive layer exists.

For a small fixed scaling \(\delta>0\), use the optimization of
\eqref{eq:E3} with
\begin{equation}\label{eq:G1}
 B_*=\delta B_0,\qquad q_*=\delta/4,\qquad
 q_{\rm bar}=\delta/16,\qquad
 \mathcal L_\delta(w)=A_\theta(w)+\delta^{-1}(q_*-q(w)).
\end{equation}
Only finitely many fixed scalings will be needed before taking
the sequence limit.  We work with objective bounds strictly
below \(1/8\).  At every such optimizer,
\[
                         q(w)>\delta/8,\qquad A_\theta(w)<1,
\]
because both summands in \eqref{eq:G1} are nonnegative.
At the cheap parameters \(t=0\), the starting optimum is at
most \(\delta\kappa\): use \(\delta w_c\), whose barrier is
at least \(\delta b/2\), so the cap is attained.

Consider a layer with scale \(S\), and let \(s\) be the number
of axes up to and including this layer.  Inductively, the old
probes, from cheap and preceding layers, are actual angular
functions of degree and \(v\)-cost at most \(C W_j\), and their
Gaussian parameters are integral, satisfy
\(|t_j|\le C\delta W_j\), and are zero on the cheap axes.
Suppose the old minimum of \(\mathcal L_\delta\) is less than \(1/8\).
Constants below may depend on previous size bounds and fixed
width comparabilities, but not on the small fixed \(\delta\).
Choose an optimizer \(w\) with the tie rule from
Section~\ref{c:optimization-section}.  The two-sided estimate
\eqref{eq:F4} applies to it.

For \(J\ge1\), let \(E_J\subset K_w\) be the span of the
homogeneous \(w\)-initials of elements of
\(V_{\lceil JS\rceil_n}\) with value at most \(K_2\delta JS\).
Here \(\lceil x\rceil_n\) denotes the least nonnegative multiple
of \(n\) at least \(x\), and \(K_2\) is a sufficiently large
fixed constant.  Nesting of the systems makes these spaces
nested in \(J\).

\begin{lemma}[Layer dimension]\label{c:layer-dimension}
For every fixed \(J\ge1\), on sufficiently late members,
\begin{equation}\label{eq:G2}
                   \dim E_J\asymp
                      J^s\prod_{j\le s}\frac{S}{W_j}.
\end{equation}
The comparison constants are independent of \(J\) and
\(\delta\); the required lateness may depend on both.
\end{lemma}

\begin{proof}
At the tested degree, the barrier \(q>\delta/8\) sends costs
\(b>C'JS\), for a sufficiently large fixed \(C'\), to values
\(w>K_2\delta JS\).  Equation~\eqref{eq:C8} bounds the quotient
through the former cost by a constant times the right side
of \eqref{eq:G2}, since fixed \(JS\) is negligible compared
with widths in subsequent layers.  This proves the upper
bound.  In particular the quotient by \(w>K_2\delta JS\),
with zero included in the high part, is finite-dimensional.
Its dimension equals \(\dim E_J\), by taking initials on its
successive value slices.  Distinct occurring values already
give independent quotient classes.

For the lower bound, the original independent probe boxes
following \eqref{eq:C8} and \eqref{eq:D3} give products with
exponents
\[
               0\le a_j\le\left\lfloor c'JS/W_j\right\rfloor
                         \qquad(j\le s),
\]
for a sufficiently small fixed \(c'>0\).
After nesting, they lie in the tested physical degree.
They are independent, and their \(v\)-Gaussianity bounds the
cost of every nonzero linear combination by \(C''JS\).
Equation~\eqref{eq:F4} bounds its absolute \(w\)-value by
\(C_6C''\delta JS\).  Taking \(K_2\) above this constant
makes their span inject into the quotient.  Their number
is bounded below by the required box product.
\end{proof}

\begin{proposition}[Rank increase in a layer]\label{c:rank-increase}
There is a bounded \(J\), independent of the small fixed
\(\delta\), such that the old probe initials extend to \(s\)
algebraically independent homogeneous initials in \(E_J\).
Their actual lifts can be chosen with degree and cost
\(O(S)\), and absolute \(w\)-value \(O(\delta S)\).
\end{proposition}

\begin{proof}
Choose a fixed \(J_0\ge1\) whose space contains the old
probe initials; their widths are \(o(S)\).
Suppose the transcendence degrees of the fields generated
by \(E_{J_0}\) and \(E_N\) are equal to \(a<s\), where \(N\)
is a sufficiently large fixed integer to be chosen in terms
of \(J_0\).  The dimension of \(E_N\) grows as \(N^s\) times
the layer box factor in \eqref{eq:G2}.  We will encode its elements
by Taylor orders in \(a\) variables and scale their order polytope
by \(1/N\).  This loses only a factor \(N^a\).  Trace will put
the resulting independent outputs into a degree and value budget
independent of \(N\), contradicting the upper section count when
\(a<s\).

Choose a transcendence basis \(g=(g_j)_{j\le a}\)
from the homogeneous generators of \(E_{J_0}\), extending
the old probes, with their actual lifts in that budget.
If \(a=0\), every element of \(E_N\) is algebraic over the
algebraically closed constant field, so \(\dim E_N\le1\).
This contradicts \eqref{eq:G2} for large \(N\), since
\(S/W_j\) is bounded below for \(j\le s\).
Assume henceforth \(a>0\).
The finite space \(E_N\) lies in a finite separable extension
of \(k(g)\).

Extend constants independently to \(k'=\overline{k(\xi)}\),
with \(\xi_1,\ldots,\xi_a\) independent, and use the generic
component at \(g=\xi\) from the Taylor construction
\eqref{eq:E9}.  Put \(u'_j=g_j-\xi_j\).
Every nonzero pre-extension element of the finite algebraic
field is a unit at the generic point of this component.
Its Taylor expansion has coefficients in \(k'\).
Indeed take its separable polynomial over \(k(g)\).
At this generic specialization the coefficients are regular,
the chosen root remains simple, and its residue lies in
the algebraically closed field \(k'\).  Simple-root lifting
gives the unique formal solution in \(k'[[u']]\).
This also describes all \(k'\)-linear combinations.

This operation keeps full Taylor series, not merely residue
values.  Geometric integrality preserves linear independence
under the independent constant extension, and the local
completion is injective on the functions under consideration.
For example the independent functions \(1,g_j\) expand as
\(1,\xi_j+u'_j\), rather than being replaced by \(1,\xi_j\).
Thus the extended \(E_N\) retains its dimension and has
exactly \(\dim_k E_N\) distinct lexicographic Taylor orders.
Each leading-order slice has dimension one over \(k'\),
since its leading coefficient is a scalar; elimination then
gives the equality of counts.  These orders lie in
\(\Z_{\ge0}^a\) and contain \(0\) and all standard basis
vectors, using \(1,u'_j\).

Let \(\mathcal P\) be their convex hull.  It is full-dimensional
and
\[
                         \vol(\mathcal P)\ge c_d\dim E_N.
\]
To verify this bound even when many points lie on faces,
triangulate using all the given lattice points as vertices,
inserting them by successive subdivisions.  Every such
vertex lies in a full simplex, so there are at least
\(\dim E_N/(a+1)\) full simplices.  Each lattice simplex
has volume at least \(1/a!\).

Translation averaging gives \(z'\in[0,1]^a\) such that
\(z'+\mathcal P/N\) contains at least
\(\vol(\mathcal P)/N^a\) lattice points \(b'\).
They are nonnegative.  For each occurring order \(\alpha_h\)
choose an element \(h'\) of the extended \(E_N\) with that
order and a nonzero scalar leading coefficient.
Retain their finite Taylor leading data and their expressions
in homogeneous initials with actual lifts, together with
those of \(g\), in the spread for
\eqref{eq:E7}--\eqref{eq:E9}.

We require shifts for whatever Frobenius power
\(P_0=\mathfrak p^{\ell_0}\), \(\ell_0\ge1\), that construction later supplies.
Write
\[
        N(b'-z')=\sum_h\lambda_h\alpha_h,\qquad
        \lambda_h\ge0,\qquad\sum_h\lambda_h=1.
\]
Use \(h'\) to exponent \(k_h=\lfloor P_0\lambda_h/N\rfloor\).
The remaining Taylor order is
\[
 \begin{split}
  P_0b'-\sum_h k_h\alpha_h
      &=P_0z'+\sum_h
          \left\{\frac{P_0\lambda_h}{N}\right\}\alpha_h.
 \end{split}
\]
It is integral and componentwise nonnegative.
Its \(j\)-th coordinate is at most
\(P_0+\sum_h(\alpha_h)_j\), hence at most \(2P_0\)
once \(\mathfrak p\) exceeds the finitely many coordinate sums.
This one choice works for every \(\ell_0\ge1\).
Fill the remainder by powers of \(u'_j\).
The resulting shift has exact Taylor order \(P_0b'\),
a nonzero scalar leading coefficient, total \(E_N\)
multiplicity at most \(P_0/N\), and at most \(2P_0\)
filler factors per coordinate.

Use Corollary~\ref{c:strict-budget} and
Proposition~\ref{c:trace-budget}
over the probes \(g\), with \(\kappa=1/\delta\), at the
same old optimizer.  Only the old probes are optimization
constraints; enlarging the list for the relative field
does not alter this optimization.
Apply every shift to the same surviving relative input
and the same trace power from \eqref{eq:E7}.
The full initial-field outputs have the distinct Taylor
orders \(b'\), so are independent over the enlarged
perfect constants.  Their common relative leading
coefficient may be nonconstant in the residue field;
it is nonzero, which is all that distinct-order
independence needs.  Semilinearity only takes the
\(P_0\)-th roots of scalar coefficients.

Expand the inputs in homogeneous angular initials and
discard summands with zero initial-field trace.
Each remaining summand lifts by equality in
\eqref{eq:E8} to an actual output initial.
Its physical degree and value satisfy
\[
       m'\le C_7(1+J_0)S,\qquad
       w(F')\le C_7(1+J_0)\delta S.
\]
These estimates are termwise.  Each \(E_N\) generator has
degree at most \(\lceil NS\rceil_n\) and value at most
\(K_2\delta NS\), and there are at most \(P_0/N\)
occurrences.  The filler factors and the flags of
\eqref{eq:E9} use the \(E_{J_0}\) budget with at most
\(3P_0\) occurrences per axis.  Expand \(g_j-\xi_j\)
as a probe plus a scalar before estimating; both terms
satisfy that budget.  Likewise use the fixed bounded
homogeneous expressions for each \(h'\), so no
cancellation is needed to improve an upper estimate.

The active terms of \eqref{eq:E10}, after division by \(P_0\),
cost only \(O(1/\delta)\) in degree and \(O(1)\) in value.
Indeed \(m_i/b_i\) is bounded and
\(w(F_i)=qb_i-\delta B_0m_i\).
The tiny factors and active-exponent approximation errors
can be made arbitrarily small.  Since \(\delta\) is fixed
before \(S\to\infty\), these contributions are negligible
in the displayed budgets.  Apply \eqref{eq:E8} with
equality, using \(A_\theta(w)\ge0\), and its degree bound.

Consequently the independent initial-field outputs belong
to the span of actual output initials in those two budgets.
Before choosing the characteristic, pad to one nested
degree dominating \(C_7(1+J_0)S\), retain the cost-dimension
bound through \(C'(1+J_0)S\) at \(v\), and retain the
inclusion of the larger-cost subspace into values strictly
above the required cutoff at \(w\).
The barrier \(q>\delta/8\) gives this inclusion with margin
for sufficiently large fixed \(C'\).  Lemma~\ref{c:trace-finite-data}, with the budget
accounting of Proposition~\ref{c:trace-budget}, therefore
bounds that output span in reduction by
\[
                 C_8(1+J_0)^d\prod_{j\le s}S/W_j.
\]
It applies to all coefficient vectors of this predetermined
space, not to a list of outputs chosen after \(P_0\).
On the other hand, the number of distinct orders \(b'\)
is at least
\[
                       cN^{s-a}\prod_{j\le s}S/W_j
\]
by \eqref{eq:G2}.  Choose \(N\) large in terms of \(J_0\)
and the preceding constants, independently of \(\delta\).
These bounds contradict each other.

Thus rank increases from \(E_{J_0}\) to \(E_{N(J_0)}\)
unless it has reached \(s\).  Iterate, setting the next
\(J_0\) equal to that \(N\), at most \(s\) times, all at
the same optimizer.  Choose the new probes among the
homogeneous generators, extending the old list.
Their degree is \(O(S)\); the barrier bounds their
cost by \(CS\) and their value below by
\(-\delta B_0 O(S)\), while the defining cutoff bounds
it above by \(C\delta S\).  Their initials are independent,
so \(w\) is Gaussian on the enlarged list.  This proves
all assertions.
\end{proof}

\begin{proposition}[Slopes after a rank increase]
\label{c:layer-slopes}
For the enlarged optimization with \(s\) probes, on
full-dimensional generic cells with
\(|t_j|\le C\delta W_j\) and minimum \(<1/8\),
\begin{equation}\label{eq:G3}
       \left|\partial_{t_j}\min\mathcal L_\delta\right|
                            \le C_9\,S/W_j
                            \qquad(j\le s).
\end{equation}
The same assertion holds in slightly enlarged fixed
parameter boxes and at each of the later fixed scalings.
In particular the new-coordinate slopes are uniformly
bounded independently of the small fixed \(\delta\).
\end{proposition}

\begin{proof}
We repeat the slope mechanism from the proof of
Proposition~\ref{c:cheap-collapse}, using rectangular probe boxes
whose side lengths reflect the widths
\(W_j\).  We retain its central trace equality and nearby lower
inequality.  Here even cheap parameters may vary, so the lower section
barrier supplies the upper dimension count without using
\eqref{eq:F4}.
Choose a generic parameter \(t\)
whose coordinates are \(\Q\)-linearly independent, and
an optimizer with the tie rule.  Use its affine
continuation with fixed stratum cone and active mask
from Proposition~\ref{c:optimization}.
Apply Corollary~\ref{c:strict-budget} and
Proposition~\ref{c:trace-budget} over
\(\overline y\), and lift a surviving initial trace
by \eqref{eq:E8}.  Its actual output \(F'\), in the
reduction, satisfies
\[
                         m'\le C'S,\qquad
                         w(F')\le C'\delta S.
\]
The active contributions cost \(O(1/\delta)\) and
\(O(1)\) in the two budgets.  The flag powers divided
by \(P_0\) cost \(O(S)\) and \(O(\delta S)\), by
the probe and parameter bounds.  Choose tiny factors
and approximation errors after knowing this continuation,
including their Lipschitz errors.  For fixed \(\delta\),
the claimed bounds follow on late members.

Fix a sufficiently large integer \(N\), and consider
\[
       (F')^h y^\alpha,\qquad
       0\le h<N,\qquad
       0\le\alpha_j<L'_j:=\lceil NS/W_j\rceil.
\]
They lie in one nested degree \(O(NS)\) and have
upper value \(O(\delta NS)\).  After enlarging the
bounds, these statements hold nearby on the same
continuation, by continuity after \(F'\) is chosen.
The quotient dimension through that value is at most
\(C''\prod_{j\le s}L'_j\).
Indeed, before applying \eqref{eq:E7}, retain
\eqref{eq:C8} in that predetermined degree and
the inclusion of costs \(b>C'''NS\) into larger
\(w\)-values, with strict margin.  The barrier
\(q>\delta/8\) supplies this inclusion, and fixed
\(NS\) is smaller than all subsequent widths.
The constants do not depend on \(N,\delta\).
Finite-space transfer with margin gives the same
bound locally on the reduced continuation.

There are therefore at most \(C'''\prod_{j\le s}L'_j\)
distinct product values.  At independent \(t\), this
forces a unique expression \(w(F')=\gamma\cdot t\)
with \(\gamma\in\Q^s\): otherwise all \(N\) layers
have different values, contradicting the count for
large \(N\).  We record the quantitative consequence.
Let \(D'\) be the order of \(\gamma\) in
\(\Q^s/\Z^s\).  Modulo \(\Z^s\), the first \(N\)
translated boxes have \(\min(N,D')\) distinct classes,
each with \(\prod_jL'_j\) points.  Thus \(D'\) is
uniformly bounded.  Restricting \(h\) to multiples
of \(D'\), the union of projections to coordinate
\(j\) contains
\[
 L'_j+\lfloor (N-1)/D'\rfloor
                  \min(L'_j,D'|\gamma_j|)
\]
points, since these are translates of an integer
interval.  Over each projected point lies a full
box with \(\prod_{l\ne j}L'_l\) points in the other
coordinates.  Comparison with the total bound, for
sufficiently large uniformly chosen \(N\), gives
\[
                    |\gamma_j|\le C S/W_j.
\]
The same argument works at nearby independent
parameters on the continuation.  There are only
finitely many possible rational vectors with the
stated denominator and coordinate bounds.  Uniqueness
at the central independent parameter and continuity
therefore show that, in a neighborhood of that parameter,
\[
                    w_{\widetilde t}(F')=
                                \gamma\cdot\widetilde t
\]
with this same \(\gamma\).

Let \(G\) be the input term lifted for \(F'\).
Equation~\eqref{eq:E8} makes
\[
 w_{\widetilde t}(F')-
        \frac{w_{\widetilde t}(G)}{P_0}
       +(1-1/P_0)A_\theta(w_{\widetilde t})
\]
nonnegative nearby and zero centrally.  In the active
part of \(w_{\widetilde t}(G)/P_0\), each
\(w_{\widetilde t}(F_i)\) varies by \(b_i\,dq\).
Equation~\eqref{eq:E10} says that the resulting
variation is \((1-1/P_0)dq/\delta\), up to an
arbitrarily small Lipschitz slope error.  If the cap
is active, \(dq=0\).  Tiny systems give arbitrarily
small further slope errors, while the flag powers
in \eqref{eq:E9} have bounded slopes.  All these
formulas use the chosen characteristic-zero continuation
and its chart valuations in reduction; finite input
data are retained and errors are chosen after that
continuation is known, as in the proof of \eqref{eq:F2}.
The displayed nonnegative function has a local minimum
at the central point.  Using both signs of variation,
the formula for \(w_{\widetilde t}(F')\), and
\eqref{eq:G1}, bounds the objective slope by
\(C S/W_j\).  Here bounded extra slopes are absorbed
because \(S/W_j\) is bounded below.  This proves
\eqref{eq:G3} on the generic cells.

Openness, continuity, and local polyhedrality on the
strict sublevel set extend this bound to changes along
segments varying only the new coordinates, even when
the old coordinates are fixed at nongeneric values.
Equivalently, approximate such a segment by generic
ones, apply the bound piecewise on the finitely many
local cells, and pass to the limit using compact
sublevels.  The assertions concerning enlarged boxes
and later fixed scalings follow from the same estimates.
\end{proof}

\begin{proposition}[Integral rounding]\label{c:integral-rounding}
For every absolute tolerance \(\sigma>0\), a rank
increase followed by a bounded integral scaling and
rounding of the new parameters changes a small objective
bound according to
\begin{equation}\label{eq:G4}
 e_{\mathrm{new}}=K(\sigma,d)e_{\mathrm{old}}
                                    +C_{10}\sigma.
\end{equation}
Old parameters remain integral and cheap parameters
remain zero.  The constants for probe and parameter
sizes remain uniform at every later fixed scaling.
\end{proposition}

\begin{proof}
Simultaneous Dirichlet approximation gives a positive
integer \(u\le K(\sigma,d)\) for which every new
parameter of \(ut\) is within \(\sigma\) of an integer.
For example, subdivide the unit cube into boxes of
side at most \(\sigma\), apply the pigeonhole principle
to the fractional parts of finitely many multiples
of the new parameter vector, and subtract two multiples
in the same box.  Multiply \(w,t,\delta\) by this
common \(u\).  Integrality of the old parameters is
preserved.  Both \(q\) and \(q_*\) scale by \(u\),
so the nonnegative penalty in \eqref{eq:G1} is unchanged;
discrepancy scales by \(u\).  Thus the starting objective
for the newly scaled enlarged optimization is at most
\(u e_{\mathrm{old}}\).

Move only the new parameters to the chosen integers
along a line segment.  Integrating \eqref{eq:G3} in
the new coordinates adds at most \(C_{10}\sigma\).
Choose tolerances so that the full resulting bound
is strictly less than \(1/8\).  If continuation had
a first limiting endpoint, compact sublevels would
supply a feasible optimizer there with this same
strict bound, and local continuation would extend
the segment.  Hence the whole rounding is possible.
At fixed scaling and tolerance, the rounding errors
are negligible relative to \(\delta W_j\) in the
new layer on sufficiently late members.  The path
therefore stays within slightly enlarged parameter
boxes, and the stated size bounds persist under
subsequent scaling.
\end{proof}

\begin{proposition}[Bounded divisorial index]\label{c:divisorial-index}
At a final optimizer with \(d\) independent probe
initials and integral probe parameters, there is a
uniformly bounded positive integer \(e\) such that
\(ew\) is a positive integral multiple of a
prime-divisor valuation over \(X\).
\end{proposition}

\begin{proof}
First bound the image of the value group of \(w\)
in \(\R/\Z\).  Choose any finite list of functions
whose values represent distinct classes.  Multiplying
them all by a common denominator puts them in \(V_0\)
and translates all classes by the same amount, so
they remain distinct.  Multiply each function by
all probe monomials in the boxes
\[
                    0\le\alpha_j\le\lfloor R/W_j\rfloor,
                    \qquad 1\le j\le d.
\]
Here \(R\to\infty\) on this one fixed member of the
sequence.  The resulting initials are linearly
independent: between different functions their values
are in disjoint classes modulo \(\Z\), since all
probe values are integral; within one box Gaussianity
of the probes gives independence.

At a common nested degree \(O(R)\) all these products
have upper \(w\)-value \(O(\delta R)\).  The constants
are uniform by the size bounds, once \(R\) is also
large enough for the finitely many chosen function
degrees and values.  The barrier \(q>\delta/8\) and
\eqref{eq:C8} bound the dimension through this cutoff
by
\[
                      C_{11}\prod_{j=1}^d(1+R/W_j).
\]
Comparing with the number of independent products
and letting \(R\) tend to infinity bounds the length
of the finite list uniformly.  Thus the image of the
entire value group modulo \(\Z\) is a finite subgroup
of uniformly bounded order.  Its exponent is a uniformly
bounded positive integer \(e\), and \(ew\) is
integer-valued.

Equation~\eqref{eq:F3} makes \(w\) nonzero on the
degree-zero field.  It is an Abhyankar valuation
with center on \(X\), as supplied throughout by
monomialization.  A nonzero discrete rank-one
Abhyankar valuation is divisorial: on a monomial
chart, the rational weight ray is extracted by a
toric subdivision, and the corresponding primitive
ray defines a prime divisor.  Since \(ew\) is
integer-valued, it is a positive integral multiple
of that prime-divisor valuation.  This proves the
claim.
\end{proof}

\phantomsection\label{c:choice-order}
\begin{proof}[Completion of Proposition~\ref{prop:birational-threshold}]
We now specify the order of all choices.  The rank,
probe-size, parameter-size, and slope constants in
\eqref{eq:G2}--\eqref{eq:G3} can be bounded successively
forward through the at most \(d\) layers.  These
bounds are independent of the sufficiently small
fixed scalings and rounding tolerances: each step
uses only the preceding uniform size bounds, with
fixed enlargements of their constants.  Fix these
constants first, including a final bound \(C_{12}\)
for the integer \(e\) in Proposition~\ref{c:divisorial-index}.

Prescribe objective tolerances backwards through
\eqref{eq:G4}.  At each step choose \(\sigma\) first
so that \(C_{10}\sigma\) is smaller than the desired
next tolerance, then choose the previous objective
bound so small that its product with \(K(\sigma,d)\)
fits the remaining allowance.  Make every intermediate
bound satisfy all strict-sublevel requirements, and
make the final bound less than \(\epsilon/(3C_{12})\).
Only after these finitely many choices choose the
initial \(\delta\) so small that \(\delta\kappa\)
is below the first starting tolerance.  Require also
that after every bounded possible product of the
integer scalings, the resulting \(\delta\) remains
small and satisfies
\[
                         \delta B_0<\epsilon/(3C_{12}).
\]
Finally pass to sufficiently late members of the
sequence.  All estimates involve fixed scalings;
only finitely many possible integer scaling products
have to be accommodated.

Cheap collapse, rank extension, and integral rounding
then give a final optimizer with
\[
 A_\theta(w)<\epsilon/(3C_{12}),
 \qquad
 A_X(w)\le A_\theta(w)+\delta B_0.
\]
The second inequality is \eqref{eq:A13} and the
budget barrier.  By Proposition~\ref{c:divisorial-index},
write \(ew=a\ord_E\), where \(E\) is a prime divisor
over \(X\), \(a\) is a positive integer, and
\(e\le C_{12}\).  Consequently
\[
           A_X(E)=\frac{e}{a}A_X(w)
                    \le e A_X(w)<\epsilon.
\]
This contradicts the assumed \(\epsilon\)-log canonicity
of \(X\).  Hence the sequence violating
\eqref{eq:A1} cannot exist, proving
Proposition~\ref{prop:birational-threshold}.
\end{proof}

Theorem~\ref{thm:main} follows by
Lemma~\ref{lem:threshold-to-complement} and
Lemma~\ref{lem:field-descent}.  The resulting complement
index \(N\) is distinct from the auxiliary degree
\(n\) fixed for the cone construction.
